\section{Strategi Pengujian}
Pengujian dilakukan dengan membangun program menggunakan \texttt{make -B}, menjalankan \texttt{arion\_parser} pada input Milestone 3, lalu membandingkan output aktual dengan file output yang tersimpan pada direktori \texttt{test/milestone-3}. Command yang digunakan adalah sebagai berikut.

\begin{tcolorbox}[title=Command Pengujian,colback=gray!5,colframe=black!60,boxrule=0.5pt,arc=2mm]
  \begin{verbatim}
make -B
for i in 1 2 3 4 5 6 7 8; do
  ./arion_parser test/milestone-3/input-$i.txt output-$i.txt
  diff -q output-$i.txt test/milestone-3/output-$i.txt
done
\end{verbatim}
\end{tcolorbox}

Seluruh test case yang digunakan berhasil menghasilkan exit code 0 dan output yang konsisten dengan expected output. Laporan ini menampilkan lima kasus utama yang mewakili deklarasi, ekspresi, control flow, structured type, dan subprogram.

\begin{longtable}{|c|p{4.8cm}|p{7.8cm}|}
  \caption{Ringkasan Pengujian Milestone 3} \label{tab:ringkasan-m3}                                                                                                       \\
  \hline
  \textbf{No} & \textbf{File Input}                    & \textbf{Fokus Pengujian}                                                                                          \\
  \hline
  \endfirsthead
  \hline
  \textbf{No} & \textbf{File Input}                    & \textbf{Fokus Pengujian}                                                                                          \\
  \hline
  \endhead
  1           & \texttt{test/milestone-3/input -1.txt} & Deklarasi variabel, assignment, ekspresi aritmatika, procedure call \texttt{writeln}, \texttt{tab}/\texttt{btab}. \\
  \hline
  2           & \texttt{test/milestone-3/input -2.txt} & Operator aritmatika, \texttt{div}, \texttt{mod}, unary minus, \texttt{not}, dan kondisi Boolean.                  \\
  \hline
  3           & \texttt{test/milestone-3/input -3.txt} & Control flow: \texttt{while}, \texttt{for downto}, dan \texttt{repeat-until}.                                     \\
  \hline
  4           & \texttt{test/milestone-3/input -4.txt} & Structured type: record, array, field access, array access, dan case statement.                                   \\
  \hline
  5           & \texttt{test/milestone-3/input -5.txt} & Procedure, function, parameter, local scope, recursive function call, dan function return.                        \\
  \hline
\end{longtable}

\section{Detail Pengujian}

\subsection{Kasus 1: Program Dasar}
\textbf{Input}\par
\lstinputlisting[basicstyle=\ttfamily\scriptsize,breaklines=true,frame=single,numbers=none]{../../test/milestone-3/input-1.txt}

\textbf{Cuplikan Output}\par
\begin{lstlisting}[basicstyle=\ttfamily\scriptsize,breaklines=true,frame=single,numbers=none]
=== tab ===
 42     Hello    program    10    0    1    0    0    41
 43         a   variable     0    0    1    0    0    42
 44         b   variable     0    0    1    0    1    43

=== btab ===
  0    44     0     0     2
  1     0     0     0     0

=== Decorated AST ===
Program('Hello') -> type:Unknown, lev:0, block:0
  VarDecl('a') -> tab_index:43, type:Integer
  VarDecl('b') -> tab_index:44, type:Integer
  Block -> type:Unknown, lev:1, block:1
    Assign
    ProcCall('writeln')
\end{lstlisting}

Kasus ini menunjukkan bahwa program dasar dapat dianalisis secara semantik. Identifier program dicatat sebagai \texttt{program}, variabel \texttt{a} dan \texttt{b} masuk ke \texttt{tab}, block utama masuk ke \texttt{btab}, dan decorated AST memiliki anotasi tipe serta block.

\begin{figure}[H]
  \centering
  \includegraphics[width=0.95\linewidth]{04-eksperimen/tc1.png}
  \caption{Bukti input dan output pengujian Kasus 1.}
\end{figure}

\subsection{Kasus 2: Ekspresi dan Boolean}
\textbf{Input}\par
\lstinputlisting[basicstyle=\ttfamily\scriptsize,breaklines=true,frame=single,numbers=none]{../../test/milestone-3/input-2.txt}

\textbf{Cuplikan Output}\par
\begin{lstlisting}[basicstyle=\ttfamily\scriptsize,breaklines=true,frame=single,numbers=none]
Program('Expressions') -> type:Unknown, lev:0, block:0
  VarDecl('x') -> type:Integer
  VarDecl('y') -> type:Integer
  VarDecl('z') -> type:Integer
  Block -> lev:1, block:1
    Assign x := (10 + 5) - 3 * 2
    Assign y := x div 2 + 1
    Assign z := -x mod 3
    If -> condition type:Boolean
      ProcCall('writeln')
\end{lstlisting}

Kasus ini menguji inferensi tipe untuk operator aritmatika, integer division, modulo, unary minus, dan operator \texttt{not}. Kondisi pada \texttt{if} berhasil dianotasi sebagai Boolean.

\begin{figure}[H]
  \centering
  \includegraphics[width=0.95\linewidth]{04-eksperimen/tc2.png}
  \caption{Bukti input dan output pengujian Kasus 2.}
\end{figure}

\subsection{Kasus 3: Control Flow}
\textbf{Input}\par
\lstinputlisting[basicstyle=\ttfamily\scriptsize,breaklines=true,frame=single,numbers=none]{../../test/milestone-3/input-3.txt}

\textbf{Cuplikan Output}\par
\begin{lstlisting}[basicstyle=\ttfamily\scriptsize,breaklines=true,frame=single,numbers=none]
Program('ControlFlow') -> type:Unknown, lev:0, block:0
  VarDecl('i') -> type:Integer
  VarDecl('sum') -> type:Integer
  Block -> lev:1, block:1
    While -> condition type:Boolean
      Assign sum := sum + 1
    For('i', downto)
      Assign sum := sum - i
    Repeat
      Assign sum := sum + 5
      Assign i := i + 1
      condition type:Boolean
\end{lstlisting}

Kasus ini memvalidasi bahwa semantic analyzer mampu menangani \texttt{while}, \texttt{for}, dan \texttt{repeat-until}. Ekspresi kondisi pada loop diuji agar menghasilkan tipe Boolean.

\begin{figure}[H]
  \centering
  \includegraphics[width=0.95\linewidth]{04-eksperimen/tc3.png}
  \caption{Bukti input dan output pengujian Kasus 3.}
\end{figure}

\subsection{Kasus 4: Structured Type}
\textbf{Input}\par
\lstinputlisting[basicstyle=\ttfamily\scriptsize,breaklines=true,frame=single,numbers=none]{../../test/milestone-3/input-4.txt}

\textbf{Cuplikan Output}\par
\begin{lstlisting}[basicstyle=\ttfamily\scriptsize,breaklines=true,frame=single,numbers=none]
=== atab ===
idx  xtyp  etyp  eref  low  high  elsz  size
  0     7     0     0     1     5     1     5

Program('StructuredTypes') -> type:Unknown
  TypeDecl('Point') -> type:Record
    RecordType
      VarDecl('x')
      VarDecl('y')
  VarDecl('p') -> type:Record
  VarDecl('scores') -> type:Array
  VarDecl('grade') -> type:Char
  Block
    FieldAccess(.x)
    FieldAccess(.y)
    ArrayAccess scores[1]
    Case
\end{lstlisting}

Kasus ini menguji record, array, akses field, akses indeks array, serta case statement. Informasi array masuk ke \texttt{atab}, sedangkan field record disimpan melalui block record.

\begin{figure}[H]
  \centering
  \includegraphics[width=0.95\linewidth]{04-eksperimen/tc4.png}
  \caption{Bukti input dan output pengujian Kasus 4.}
\end{figure}

\subsection{Kasus 5: Subprogram}
\textbf{Input}\par
\lstinputlisting[basicstyle=\ttfamily\scriptsize,breaklines=true,frame=single,numbers=none]{../../test/milestone-3/input-5.txt}

\textbf{Cuplikan Output}\par
\begin{lstlisting}[basicstyle=\ttfamily\scriptsize,breaklines=true,frame=single,numbers=none]
Program('Subprograms') -> type:Unknown
  VarDecl('result') -> type:Integer
  ProcedureDecl('PrintTwice') -> block:1
    Param('msg': string)
    Param('count': integer)
    VarDecl('i') -> type:Integer
    For('i', to)
      ProcCall('writeln')
  FunctionDecl('Factorial', returns: 'integer') -> type:Integer, block:2
    Param('n': integer)
    If
      Assign Factorial := 1
      Assign Factorial := n * Factorial(n - 1)
  Block -> block:3
    ProcCall('PrintTwice')
    FuncCall('Factorial') -> type:Integer
    ProcCall('writeln')
\end{lstlisting}

Kasus ini menunjukkan bahwa procedure dan function dicatat dalam symbol table, parameter diproses pada scope subprogram, pemanggilan recursive function dikenali sebagai \texttt{FuncCall}, dan function \texttt{Factorial} memiliki assignment ke nama function sebagai return value.

\begin{figure}[H]
  \centering
  \includegraphics[width=0.95\linewidth]{04-eksperimen/tc5.png}
  \caption{Bukti input dan output pengujian Kasus 5.}
\end{figure}

\section{Pengujian Error Semantic}
Selain kasus valid, program juga diuji secara manual terhadap input yang mengandung kesalahan semantik. Contoh kesalahan yang berhasil dideteksi adalah undefined identifier, assignment dengan tipe tidak kompatibel, kondisi \texttt{if} non-Boolean, redeklarasi identifier, index array dengan tipe salah, dan pemanggilan by-reference dengan literal.

\begin{lstlisting}[basicstyle=\ttfamily\scriptsize,breaklines=true,frame=single,numbers=none]
Input:
program BadType;
var x: integer; s: string;
begin
  x := s;
end.

Output:
=== Semantic Errors ===
- Assignment types are incompatible
\end{lstlisting}
